\documentclass[../../main.tex]{subfiles}
\begin{document}

% SOURCE p=225; printed=210; object=ch06-notes-riesz-thorin-history; region=upper
\textbf{关于\cref{section:ch06-lp-interpolation}。}M.~Riesz（F.~Riesz 的弟弟）\cite{folland-bib-118}最先在 $p_j\leqslant q_j$（$j=0,1$）的假设下证明了 Riesz–Thorin 定理。一般情形的证明，以及使用三线引理的想法，来自 Thorin\cite{folland-bib-149}。

% SOURCE p=225; printed=210; object=ch06-notes-stein-family; region=upper-middle
E.~M.~Stein 证明了 Riesz–Thorin 定理的一个强有力的推广。它处理算子族 $\{T_z:0\leqslant\operatorname{Re}z\leqslant1\}$：粗略地说，这些算子全纯地依赖于 $z$，并在 $|\operatorname{Im}z|\to\infty$ 时满足某些温和的增长条件。定理断言，若在 $\operatorname{Re}z=j$ 上 $T_z$ 从 $L^{p_j}$ 到 $L^{q_j}$ 有界（$j=0,1$），则在 $\operatorname{Re}z=t$ 上它从 $L^{p_t}$ 到 $L^{q_t}$ 有界（$0<t<1$），其中 $p_t,q_t$ 按 Riesz–Thorin 定理的方式定义。精确陈述和证明可见 Bennett 与 Sharpley\cite[第4.3节]{folland-bib-15}、Stein 与 Weiss\cite[第V.4节]{folland-bib-142}或 Zygmund\cite[第XII.1节]{folland-bib-167}。这些思想的进一步推广见 Coifman 等\cite{folland-bib-28}。
% SUPPLEMENT: id=v14-ch06-stein-precise-version; parent=ch06-notes-stein-family; kind=statement+proof+motivation+hypotheses; origin=independent-derivation
\begin{proposition}[全纯算子族插值的一个精确版本（重构新增）]\label{proposition:ch06-stein-uniform-edges}
设 $X,Y$ 为 $\sigma$ 有限测度空间，$\Sigma_X,\Sigma_Y$ 分别为有限测度集上支承的简单函数空间。设每个 $T_z$ 为从 $\Sigma_X$ 到 $Y$ 上可测函数的线性映射，且对所有 $f\in\Sigma_X,g\in\Sigma_Y$，配对
\begin{align*}
F_{f,g}(z)=\int_Y(T_zf)g\,\mathrm d\nu
\end{align*}
绝对收敛、在闭带连续并在内部全纯。假设对每一对 $f,g$ 都存在 $A>0$ 和 $0\leqslant b<\pi$，使全带上
\begin{align*}
|F_{f,g}(x+iy)|\leqslant\exp(Ae^{b|y|}).
\end{align*}
设 $1\leqslant p_j,q_j\leqslant\infty$，且存在与 $y,f$ 无关的有限常数 $M_j\geqslant0$，满足
\begin{align*}
\|T_{j+iy}f\|_{q_j}\leqslant M_j\|f\|_{p_j}
\qquad(j=0,1;\ y\in\mathbb R;\ f\in\Sigma_X).
\end{align*}
按倒数的线性插值定义 $p_t,q_t$。若 $0<t<1$ 且 $p_t<\infty$，则每个 $T_{t+is}$ 都唯一延拓为 $L^{p_t}(X)\to L^{q_t}(Y)$ 的有界线性算子，范数不超过 $M_0^{1-t}M_1^t$。
\end{proposition}
\begin{proof}
先扩充三线引理允许的增长范围。设标量函数 $F$ 连续到闭带、内部全纯，满足上述增长条件，且两条边上分别有界于正数 $M_0,M_1$。取 $b<k<\pi$，令
\begin{align*}
G_\varepsilon(z)=F(z)\exp\bigl(-\varepsilon\cos(k(z-\tfrac12))\bigr).
\end{align*}
若 $z=x+iy$，则
\begin{align*}
\operatorname{Re}\cos(k(z-\tfrac12))
=\cos(k(x-\tfrac12))\cosh(ky)
\geqslant\cos(k/2)\cosh(ky)>0.
\end{align*}
由于 $k>b$，这个衰减因子使 $G_\varepsilon$ 在全带有界；在两条边上它的模不超过原来的 $M_j$。由\cref{lemma:ch06-three-lines}，$|G_\varepsilon(t+is)|\leqslant M_0^{1-t}M_1^t$。对固定 $t,s$ 令 $\varepsilon\downarrow0$ 得同样的 $F$ 界。边界常数为零时先以 $M_j+\delta$ 代替，再令 $\delta\downarrow0$。

固定 $t,s$，取非零 $f\in\Sigma_X,g\in\Sigma_Y$，归一化为 $\|f\|_{p_t}=\|g\|_{q_t'}=1$。令
\begin{align*}
\alpha(z)&=(1-z)/p_0+z/p_1,&
\beta(z)&=(1-z)/q_0+z/q_1.
\end{align*}
按不交集及非零系数写 $f=\sum c_\ell\chi_{E_\ell}$、$g=\sum d_m\chi_{H_m}$，并令
\begin{align*}
f_z&=\sum_\ell |c_\ell|^{\alpha(z)/\alpha(t)}
\operatorname{sgn}(c_\ell)\chi_{E_\ell},\\
g_z&=\sum_m |d_m|^{(1-\beta(z))/(1-\beta(t))}
\operatorname{sgn}(d_m)\chi_{H_m}.
\end{align*}
若 $\beta(t)=1$，则 $q_0=q_1=1$，此时改取 $g_z=g$。因为 $p_t<\infty$，第一式的分母总为正。计算边界上的模并积分，得到 $\|f_{j+iy}\|_{p_j}=1$；当 $p_j=\infty$ 时，各非零分块上的模为一，结论仍成立。同理 $\|g_{j+iy}\|_{q_j'}=1$，包括上述 $q_j'=\infty$ 的分支。

令 $\Phi(z)=\int (T_{z+is}f_z)g_z\,\mathrm d\nu$。展开有限和，它是常数示性函数配对 $F_{\chi_{E_\ell},\chi_{H_m}}(z+is)$ 乘以全纯的指数系数的有限和。因此它连续到闭带、内部全纯。取有限多个增长参数的最大值，再选仍小于 $\pi$ 的更大参数 $b'$，即可把有限和、固定平移 $is$ 和指数系数引起的至多 $e^{C|y|}$ 增长一并控制为 $\exp(A'e^{b'|y|})$。由边界算子界及 Hölder 不等式，$|\Phi(j+iy)|\leqslant M_j$。刚证出的三线版本给出
\begin{align*}
\left|\int (T_{t+is}f)g\,\mathrm d\nu\right|
\leqslant M_0^{1-t}M_1^t.
\end{align*}
对简单测试函数 $g$ 取上确界，由\cref{theorem:ch06-dual-test-criterion}及 $Y$ 的 $\sigma$ 有限性，得到 $\|T_{t+is}f\|_{q_t}\leqslant M_0^{1-t}M_1^t\|f\|_{p_t}$。$\Sigma_X$ 在有限指数的 $L^{p_t}$ 中稠密，而 $L^{q_t}$ 完备，故有唯一的连续线性延拓。
\end{proof}
\begin{remark}
\textbf{动机与逐条件说明（重构新增）。}这是为原文概述补出的一个明确、可直接证明的版本，不把它冒称为所引专著中最一般的陈述。全纯依赖须通过每个标量配对说明；只给算子符号加上复下标不能替代这一条件。配对的绝对收敛使有限展开和积分有定义；带内全纯与闭带连续使三线论证可用。增长指数严格小于 $\pi$，使能选择 $b<k<\pi$，同时获得向上下的衰减和全带上正的 $\cos(k(x-1/2))$。统一的边界常数控制整条竖线；原文“粗略地说”中的增长和边界条件因此都被落实。这里采用 $\sigma$ 有限性以保证对偶测试，属于一个充分的假设版本；$p_t<\infty$ 专门用于简单函数稠密的延拓步骤。$q_t=\infty$ 则可以保留，因为对偶测试及 $L^\infty$ 的完备性仍可用。令 $T_z$ 与 $z$ 无关，便回到三线法证明固定算子插值的思路。
\end{remark}
% END SUPPLEMENT: id=v14-ch06-stein-precise-version

% SOURCE p=225; printed=210; object=ch06-notes-marcinkiewicz-history; region=middle
Marcinkiewicz\cite{folland-bib-97}在 $p_j=q_j$（$j=0,1$）的情形宣布了他的插值定理。他在第二次世界大战中不幸早逝，此项工作后来由 Zygmund\cite{folland-bib-166}完成。这个定理对 $T$ 的假设还可以进一步减弱：对我们已经给出的论证稍作变通，只假设对某个常数 $C$ 有
\begin{align*}
|T(f+g)|\leqslant C(|Tf|+|Tg|),
\end{align*}
也可得相同结果。参见 Zygmund\cite{folland-bib-166}及\cite[第XII.4节]{folland-bib-167}。
% SUPPLEMENT: id=v14-ch06-marcinkiewicz-no-homogeneity; parent=ch06-notes-marcinkiewicz-history; kind=proof+motivation+hypotheses; origin=independent-derivation
\begin{remark}
\textbf{去掉正齐次性的完整论证（重构新增）。}这里的“只假设”是将\cref{theorem:ch06-marcinkiewicz}对映射的次线性要求换为上述准次可加不等式；原定理的定义域、可测性、指数条件与两个弱型界仍须保留。把常数增大后可设 $C\geqslant1$。对于范数为一的输入 $u$，原证明将 $u$ 分成 $g_A+h_A$。无须任何齐次性便有
\begin{align*}
\lambda_{Tu}(2C\alpha)
\leqslant\lambda_{Tg_A}(\alpha)+\lambda_{Th_A}(\alpha).
\end{align*}
原证明在归一化输入上，其余步骤只使用这条分布估计、两个弱端点界及对 $g_A,h_A$ 范数的计算。因此有限输出端点时，只需将层积分变量的因子 $2^q$ 换成 $(2C)^q$；有无穷输出端点时，同样将输出阈值 $2\alpha$ 换成 $2C\alpha$。包括前面逐一补出的三个无穷端点分支，都给出 $\|Tu\|_q\leqslant C B_p$。若 $p_0=p_1$，\cref{exercise:ch06-marcinkiewicz-equal-input}直接比较两个弱界的证明连准次可加性也不需要。

尚不能由 $T(Fu)=FTu$ 恢复一般输入，因为此等式已不在假设中。正确做法是：给定 $f$，若 $F=\|f\|_p>0$，另定义整张映射
\begin{align*}
S_F(u)=F^{-1}T(Fu).
\end{align*}
它满足同一个准次可加界。每个弱端点界也保持原常数，因为对有限 $q_j$，
\begin{align*}
[S_F(u)]_{q_j}=F^{-1}[T(Fu)]_{q_j}
\leqslant F^{-1}C_j\|Fu\|_{p_j}=C_j\|u\|_{p_j};
\end{align*}
当 $q_j=\infty$ 时，用本质上界范数作完全相同的计算。这里使用的是函数空间范数的齐次性，未使用 $T$ 的齐次性。将已经证明的归一化结论应用于 $S_F$ 和 $u=f/F$，便得
\begin{align*}
\|Tf\|_q=F\|S_F(f/F)\|_q\leqslant C B_p\|f\|_p.
\end{align*}
若 $F=0$，输入为定义域中的零元，任一个弱端点界已给出 $T0=0$ 几乎处处，结论亦成立。两个端点常数与 $C$ 不随 $F$ 改变，所以该论证也保留原定理在非退化输入指数区间的常数增长性质；输入指数相同的情形仍应记常数为 $B_{p,q}$。

动机是把某一次输入的大小吸收到一个新的算子中，而不是假定原算子具有缩放对称性。统一的准次可加常数只放大输出阈值；弱端点界对所有输入成立，才允许把同样的证明应用于每个 $S_F$。这说明正齐次性是方便的归一化工具，却不是这个插值结论的必要假设。
\end{remark}
% END SUPPLEMENT: id=v14-ch06-marcinkiewicz-no-homogeneity

% SOURCE p=225; printed=210; object=ch06-notes-lorentz-spaces; region=middle-lower
$L^p$ 与弱 $L^p$ 都是一个双参数函数空间族 $\{L(p,q):1\leqslant p,q\leqslant\infty\}$ 的成员。这些空间称为 Lorentz 空间，满足 $L^p=L(p,p)$ 和弱 $L^p=L(p,\infty)$。Marcinkiewicz 定理可以推广为关于 $L(p,q)$ 空间上算子插值的结果。参见 Bennett 与 Sharpley\cite[第4.4节]{folland-bib-15}或 Stein 与 Weiss\cite[第5.3节]{folland-bib-142}。
% SUPPLEMENT: id=v14-ch06-lorentz-definition; parent=ch06-notes-lorentz-spaces; kind=definition+proof+motivation+hypotheses; origin=independent-derivation
\begin{definition}[Lorentz 空间的一种重排约定（重构新增）]\label{definition:ch06-lorentz-spaces}
对 $1\leqslant p<\infty$、$1\leqslant s\leqslant\infty$，用\cref{exercise:ch06-decreasing-rearrangement}中定义的 $f^*$ 定义
\begin{align*}
\|f\|_{p,s}&=\left(\int_0^\infty[t^{1/p}f^*(t)]^s\frac{\mathrm dt}{t}\right)^{1/s}
&& (s<\infty),\\
\|f\|_{p,\infty}&=\sup_{t>0}t^{1/p}f^*(t).
\end{align*}
$L(p,s)$ 为上述量有限的可测函数按几乎处处相等所得的空间；另约定 $L(\infty,\infty)=L^\infty$。此处使用字母 $s$ 表示第二参数，以免与当前插值定理的输出指数混淆。
\end{definition}
\begin{remark}
\textbf{两个基本恒等式、端点与动机（重构新增）。}当 $s=p<\infty$ 时，由重排的等分布性质及层积分公式，
\begin{align*}
\|f\|_{p,p}^p
&=\int_0^\infty f^*(t)^p\,\mathrm dt
=p\int_0^\infty\alpha^{p-1}\lambda_{f^*}(\alpha)\,\mathrm d\alpha\\
&=p\int_0^\infty\alpha^{p-1}\lambda_f(\alpha)\,\mathrm d\alpha
=\|f\|_p^p.
\end{align*}
非负积分允许无穷值，故此处无需预先假定 $f$ 可积。\subcref{exercise:ch06-decreasing-rearrangement}{item:ch06-rearrangement-weak-norm}则直接给出 $\|f\|_{p,\infty}=[f]_p$，于是原文所说的两个空间恒等式成立。

第一个参数选择按水平集大小缩放的幂，第二个参数选择如何汇总不同尺度：积分聚合给出强空间，取上确界保留最坏尺度，给出弱空间。$s\geqslant1$ 是本段所讨论的参数范围；它并不自动保证这里写出的具体表达式满足三角不等式，尤其 $p=1,s=\infty$ 就是前面已讨论的弱 $L^1$ 量。因此不能只凭记号 $\|\cdot\|_{p,s}$ 就称其为范数。

原文在介绍整个双参数族时也写到了 $p=\infty$。无穷端点的约定必须另说：若把上面的积分公式直接以 $t^{1/\infty}=1$ 延伸到 $p=\infty,s<\infty$，所得空间只有零元。事实上，若 $f$ 不几乎处处为零，存在 $\alpha>0$ 使 $\lambda_f(\alpha)>0$，可取 $0<b<\lambda_f(\alpha)$；于是 $f^*(t)\geqslant\alpha$ 对 $0<t<b$ 成立，积分至少为 $\alpha^s\int_0^b\mathrm dt/t=\infty$。而 $\sup_{t>0}f^*(t)=\|f\|_\infty$：本质上界给出一向不等式，每个低于本质上界的正阈值又有正测度水平集，给出反向不等式。因此 $(\infty,\infty)$ 端点与原文恒等式相容。文献中其他无穷端点约定须按其定义理解，不能从有限指数公式猜出。
\end{remark}
% END SUPPLEMENT: id=v14-ch06-lorentz-definition

% SOURCE p=225; printed=210; object=ch06-notes-general-interpolation-methods; region=lower
还有许多其他 Banach 空间的“连续族”也有相应的插值定理，例如\cref{exercise:ch05-holder-banach}中讨论的 $\Lambda_\alpha$ 空间，以及\cref{section:ch09-sobolev-spaces}中讨论的 Sobolev 空间。对于一对 Banach 空间，还有两种构造“中间空间”的一般方法，分别称为“复方法”和“实方法”；它们可看作 Riesz–Thorin 定理及 Marcinkiewicz 定理的抽象形式。关于这些理论及其应用，可参见 Bergh 与 Löfström\cite{folland-bib-16}；关于实方法及其应用，还可参见 Bennett 与 Sharpley\cite{folland-bib-15}。

% SOURCE p=225; printed=210; object=ch06-notes-applications-history; region=bottom
\cref{corollary:ch06-maximal-lp}归功于 Hardy 与 Littlewood\cite{folland-bib-064}。\cref{theorem:ch06-weak-kernel}最早出现于 Folland 与 Stein\cite{folland-bib-51}，但证明的核心思想在此前数年已由 Stein 发现（例如参见 Stein\cite[第5.1节]{folland-bib-140}）；\cref{exercise:ch06-fractional-strong-weak}所讨论的特例可以追溯到 Hardy 与 Littlewood\cite{folland-bib-63}。
\end{document}
